\subsection{等度连续性的特殊判别法}

显然，等度连续（相应地，一致等度连续）集的每个子集都是等度连续（相应地，一致等度连续）的。同样，若 X 是拓扑（相应地，一致）空间而 Y 是一致空间，则 \(\mathcal{F}(X;Y)\) 的等度连续（相应地，一致等度连续）子集的每个有限并都是等度连续（相应地，一致等度连续）的。

设 \(X, X'\) 是两个拓扑（相应地，一致）空间，\(Y, Y'\) 是两个一致空间，\(f: X \to X'\) 是连续（相应地，一致连续）映射，\(g: Y \to Y'\) 是一致连续映射。

由定义立即可知，映射 \(u \rightarrow g \circ u \circ f\) （由 \(\mathcal{F}(X; Y)\) 到 \(\mathcal{F}(X'; Y')\) 中）把等度连续（相应地，一致等度连续）集变换成等度连续（相应地，一致等度连续）集。

命题 3. 设 X 是拓扑（相应地，一致）空间，\((\mathbf{Y}_{t})_{t \in I}\) 是一致空间的一个族，Y 是一个集合，且对每个 \(t \in I\) ，设 \(f_{t}\) 是 Y 到 \(\mathbf{Y}_{t}\) 中的一个映射。给 Y 赋以使所有 \(f_{t}\) 都一致连续的最粗的一致结构。\(\mathcal{F}(\mathbf{X}; \mathbf{Y})\) 的子集 H 等度连续（相应地，一致等度连续）的必要充分条件是：对每个 \(t \in I\) ，H 在映射 \(u \to f_{t} \circ u\) 下的像是 \(\mathcal{F}(\mathbf{X}; \mathbf{Y}_{t})\) 的一个等度连续（相应地，一致等度连续）子集。

这是定义 1、定义 2 以及 Y 的围域的定义的直接推论。

命题 4. 设 X、Y 是两个一致空间，H 是 X 到 Y 中的一致连续映射所成之集。设 \(\hat{X}, \hat{Y}\) 分别是 X、Y 的豪斯多夫完备化，并用 \(\tilde{H}\) 表示诸映射 \(\hat{u}: \hat{X} \to \hat{Y}\) （当 \(u\) 取遍 H 时）所成之集（第 II 章，§ 3，第 7 号，命题 15）。则 H 一致等度连续当且仅当 \(\tilde{H}\) 一致等度连续。

我们回顾，图表

\[
\begin{array}{c} \mathbf {X} \xrightarrow {u} \mathbf {Y} \\ i \downarrow \quad \quad \quad \quad \quad \quad \quad \quad \quad \quad \quad \quad \quad j \\ \hat {\mathbf {X}} \xrightarrow {\hat {u}} \hat {\mathbf {Y}} \end{array}\tag{1}
\]

是交换的，其中 \(i\) 与 \(j\) 是典型映射；此外，\(\mathbf{X}\) （相应地，\(\mathbf{Y}\) ）的一致结构是它在 \(i\) （相应地，\(j\) ）下 \(\hat{\mathbf{X}}\) （相应地，\(\hat{\mathbf{Y}}\) ）的一致结构的原像。于是 \(\mathbf{H}\) 一致等度连续当且仅当它在映射 \(u \to j \circ u\) 下的像一致等度连续（命题 3），从而我们已可限于 \(\mathbf{Y}\) 是豪斯多夫且完备的情形；此外，若 \(\tilde{\mathbf{H}}\) 一致等度连续，则 \(\mathbf{H}\) 也一致等度连续，因为它是 \(\tilde{\mathbf{H}}\) 在映射 \(\hat{u} \to \hat{u} \circ i\) 下的像；因此剩下要证的是 \(\mathbf{Y} = \hat{\mathbf{Y}}\) 时的逆命题。设 \(\mathbf{V}\) 是 \(\mathbf{Y}\) 的一个闭围域；由假设存在围域 \(\mathbf{U}\) （\(\mathbf{X}\) 的围域），使得关系 \((x, x') \in \mathbf{U}\) 与 \(u \in \mathbf{H}\) 蕴涵 \((u(x), u(x')) \in \mathbf{V}\) 。现在，若 \(\mathbf{U}'\) 是 \(\mathbf{U}\) 在 \(i \times i\) 下的像，则闭包 \(\overline{\mathbf{U}'}\) （\(\mathbf{U}'\) 在 \(\hat{\mathbf{X}} \times \hat{\mathbf{X}}\) 中的闭包）是 \(\hat{\mathbf{X}}\) 的一个围域（第 II 章，§ 3，第 7 号，命题 12）；假设表明，每当 \((z, z') \in \mathbf{U}'\) 且 \(u \in \mathbf{H}\) ，就有 \((\hat{u}(z), \hat{u}(z')) \in \mathbf{V}\) 。由于 \(\mathbf{V}\) 是闭的且 \(\hat{u}\) 连续，我们还有 \((\hat{u}(t), \hat{u}(t')) \in \mathbf{V}\) 对一切 \((t, t') \in \overline{\mathbf{U}'}\) 及一切 \(u \in \mathbf{H}\) 成立；证明完毕。

命题 5. 设 G、G' 是两个赋有各自左一致结构的拓扑群，H 是 G 到 G' 中的同态所成之集。则下列条件等价：

\begin{enumerate}
\def\labelenumi{\alph{enumi})}
\item
  H 在 G 的单位元 e 处等度连续，
\item
  H 等度连续，
\item
  H 一致等度连续。
\end{enumerate}

只需证明 a）蕴涵 c）。设 \(V'\) 是单位元 \(e'\) （\(G'\) 的单位元）的一个邻域；则由假设，存在 G 中 e 的邻域 V，使得 \(u(V) \subset V'\) 对一切 \(u \in H\) 成立；由于 H 的元素都是同态，关系 \(x^{-1}y \in V\) 蕴涵我们有

\[
(u (x)) ^ {- 1} u (y) = u (x ^ {- 1} y) \in \mathbf {V} ^ {\prime}.
\]

依据 G 与 G' 的左一致结构的围域的定义（第 III 章，§ 3，第 1 号），结论得证。

